% ============================================================
% OP3_1a_D1_blocks_v1.tex  (EXACT INSERTION BLOCKS for review;
% NOT yet implemented). Scope per GPT verdict on recon v2:
% three-notion distinction + compact computed mini-lemma +
% softened Goldstone-flatness reading + M-(alpha,beta) as
% matching question + status sentence. Excluded: admissible-class
% definition; E2/Lifshitz; direct-method/minimiser language;
% derives-claims; any claim the recorded route already solves
% OP3.1a.
% ============================================================

% ---------- ANCHOR 1 ----------
% Insert AFTER the end of the "Status of this section." paragraph
% of the dynamic-motivation subsection:
%   "The admissibility of the ordered branch as a physical realisation
%    is guaranteed by the structural property~S7 of the $\Phi$-medium,
%    while non-uniform ordered backgrounds are structurally admissible
%    by~S9 and~P6."
% and BEFORE:
%   "\subsection{Collective variables and the ordered-branch infrared EFT}"

\paragraph{Three distinct notions (recorded for OP-grad/OP3.1a).}
For the pure gradient sector $S=\int d^4X\,P(X)$ the local field
equation is $\partial_A\bigl(P_X(X)\,\partial_A\Phi\bigr)=0$; a
constant-gradient profile has $X=\mathrm{const}$ and solves it for
any constant value of $P_X$.
The recorded condition~\eqref{eq:PX_minimum} is a different
statement: it selects the magnitude $u_*^2$ as a stationary
(minimal) Euclidean action density among homogeneous gradients.
A third, again separate, question is whether the fluctuation
tensor around the ordered background matches the broken-phase EFT
tensor.

\paragraph{Computed mini-lemma (pure-$P$ fluctuation tensor).}
Writing $\Phi=\bar\Phi+\chi$ with $\partial_A\bar\Phi=u_*n_A$, so
that $\delta X=2u_*\,n\!\cdot\!\partial\chi+(\partial\chi)^2$, the
quadratic action of the pure-$P$ sector is
\[
S^{(2)}=\int d^4X\Bigl[\,P_X(X_0)\,(\partial\chi)^2
+2P_{XX}(X_0)\,u_*^2\,(n\!\cdot\!\partial\chi)^2\Bigr],
\qquad X_0=u_*^2 ,
\]
i.e.\ a fluctuation tensor
$\mathcal K^{AB}\propto P_X(X_0)\,\delta^{AB}
+2P_{XX}(X_0)\,u_*^2\,n^An^B$, with longitudinal coefficient
$P_X+2P_{XX}u_*^2$ and transverse coefficient $P_X$.
At the stationary-density point $P_X=0$, $P_{XX}>0$, the pure-$P$
sector supplies positive longitudinal stiffness but leaves the
transverse orientation/Goldstone directions flat at this order.
The recorded orientation stiffness
$\mathcal C_n\,u^2(\partial n)^2$ is therefore a natural EFT-level
supplier of the transverse kinetic term, but deriving it from the
bare bridge remains open.

\paragraph{Matching question M-($\alpha,\beta$) (a question, not
an inconsistency).}
Direct formal matching of $\mathcal K^{AB}$ to the broken-phase
kinetic tensor $\beta\,\delta^{AB}-\alpha\,n^An^B$ with
$\alpha>\beta>0$ would require $P_X(X_0)>0$ and $P_{XX}(X_0)<0$
--- not the recorded stationary-density point.
A constant-gradient profile may still solve the local pure-$P$
field equation, but the mechanism selecting its magnitude and
matching its fluctuation tensor to the broken-phase EFT tensor is
not supplied by the recorded $P_X=0$ minimum alone.
Possible bridges include mixed $P(X,\Phi)$ dependence, coupling to
the potential sector, orientation stiffness at EFT level, or an
explicit distinction between bare and EFT fluctuation tensors.
The two tensors may legitimately live at different levels; the
task recorded here is to exhibit the bridge explicitly.
OP-grad/OP3.1a sharpened; completion open.

% ---------- ANCHOR 2 (consistency echo; strike if not wanted) ----------
% In the formal OP-grad registry item, APPEND after:
%   "Dimensional analysis~\eqref{eq:u0_dimensional} is consistent,
%    but a first-principles derivation remains open."
% the sentence:
%   "A three-notion sharpening (local field equation;
%    stationary-density magnitude selection; fluctuation-tensor
%    matching) and the matching question M-($\alpha,\beta$) are
%    recorded in the dynamic-motivation subsection; completion
%    remains open."
